\documentclass{beamer}
\usetheme{Madrid}
\usecolortheme{default}

% Define OSU orange color
\definecolor{osaorange}{RGB}{220,115,38}

% Set colors
\setbeamercolor{structure}{fg=osaorange}
\setbeamercolor{title}{fg=white,bg=osaorange}
\setbeamercolor{frametitle}{fg=white,bg=osaorange}

\usepackage{graphicx}
\usepackage{amsmath}
\usepackage{booktabs}
\usepackage{hyperref}
\usepackage{tikz}
\usepackage{listings}
\usetikzlibrary{shapes,arrows}

\title{H524 Introduction to Biostatistics}
\subtitle{Week 10: More on Simple Linear Regression \& Project Presentations}
\author{Dr. John Molitor}
\institute{}
\date{Fall 2025}

\begin{document}

\frame{\titlepage}

\begin{frame}
\frametitle{Outline}
\tableofcontents
\end{frame}

\section{Review from Week 9}

\begin{frame}
\frametitle{Quick Review: Simple Linear Regression}
\textbf{Last week we covered:}
\begin{itemize}
\item Simple linear regression model: $Y = \beta_0 + \beta_1 X + \epsilon$
\item Least squares estimation: minimizing residuals
\item Coefficient interpretation
\item Correlation vs. regression
\item R-squared and model fit
\end{itemize}

\vspace{0.5cm}
\textbf{This Week:} Model diagnostics, assumptions, transformations, and presenting results

\vspace{0.3cm}
\textbf{Reading:} Pagano \& Gauvreau Chapter 18 (continued)
\end{frame}

\begin{frame}
\frametitle{The Fitted Model Recap}
\textbf{Simple Linear Regression:}
$$\hat{y} = b_0 + b_1 x$$

\vspace{0.3cm}
where:
\begin{itemize}
\item $b_1 = r \cdot \frac{s_y}{s_x}$ (slope)
\item $b_0 = \bar{y} - b_1\bar{x}$ (intercept)
\item $r$ = correlation coefficient
\end{itemize}

\vspace{0.5cm}
\textbf{Key Concepts:}
\begin{itemize}
\item \textbf{Residual:} $e_i = y_i - \hat{y}_i$ (observed - predicted)
\item \textbf{Sum of Squared Errors:} $SSE = \sum e_i^2$
\item \textbf{R-squared:} $R^2 = r^2$ (proportion of variance explained)
\end{itemize}
\end{frame}

\section{Regression Assumptions}

\begin{frame}
\frametitle{The Four Key Assumptions}
\textbf{For valid inference from linear regression:}

\vspace{0.3cm}
\begin{enumerate}
\item \textbf{Linearity:} The relationship between $X$ and $Y$ is linear
  \begin{itemize}
  \item Check: Scatterplot, residual plot
  \end{itemize}

\item \textbf{Independence:} Observations are independent
  \begin{itemize}
  \item Check: Study design, temporal patterns
  \end{itemize}

\item \textbf{Homoscedasticity:} Constant variance of residuals
  \begin{itemize}
  \item Check: Residual vs fitted plot
  \end{itemize}

\item \textbf{Normality:} Residuals are approximately normal
  \begin{itemize}
  \item Check: Q-Q plot, histogram of residuals
  \end{itemize}
\end{enumerate}

\vspace{0.3cm}
\textbf{Acronym:} LINE (Linearity, Independence, Normality, Equal variance)
\end{frame}

\begin{frame}
\frametitle{Why Assumptions Matter}
\textbf{Consequences of violated assumptions:}

\vspace{0.3cm}
\textbf{Linearity violation:}
\begin{itemize}
\item Biased estimates
\item Poor predictions
\item Misleading relationships
\end{itemize}

\vspace{0.3cm}
\textbf{Independence violation:}
\begin{itemize}
\item Underestimated standard errors
\item Invalid p-values and confidence intervals
\item False sense of precision
\end{itemize}

\vspace{0.3cm}
\textbf{Homoscedasticity violation:}
\begin{itemize}
\item Inefficient estimates
\item Invalid hypothesis tests
\item Incorrect prediction intervals
\end{itemize}

\vspace{0.3cm}
\textbf{Normality violation:}
\begin{itemize}
\item Less critical with large samples (CLT)
\item Affects confidence intervals and hypothesis tests
\end{itemize}
\end{frame}

\section{Diagnostic Plots}

\begin{frame}[fragile]
\frametitle{Residual vs Fitted Plot}
\textbf{Purpose:} Check linearity and homoscedasticity

\vspace{0.3cm}
\textbf{What to look for:}
\begin{itemize}
\item \textbf{Good:} Random scatter around horizontal line at zero
\item \textbf{Bad:} Patterns, curves, funneling
\end{itemize}

\vspace{0.5cm}
\textbf{Common patterns:}
\begin{itemize}
\item \textbf{Curved pattern:} Non-linear relationship (try transformation)
\item \textbf{Funnel shape:} Non-constant variance (heteroscedasticity)
\item \textbf{Clusters/groups:} Missing categorical variable
\item \textbf{Outliers:} Points far from the cloud
\end{itemize}

\vspace{0.3cm}
\textbf{R code:}
\begin{verbatim}
plot(fitted(model), residuals(model))
abline(h=0, lty=2)
\end{verbatim}
\end{frame}

\begin{frame}[fragile]
\frametitle{Q-Q Plot (Quantile-Quantile)}
\textbf{Purpose:} Check normality of residuals

\vspace{0.3cm}
\textbf{What to look for:}
\begin{itemize}
\item \textbf{Good:} Points fall approximately on a straight line
\item \textbf{Bad:} Systematic departures from the line
\end{itemize}

\vspace{0.5cm}
\textbf{Common patterns:}
\begin{itemize}
\item \textbf{S-shaped curve:} Heavy tails (outliers)
\item \textbf{Curve below then above:} Right skewness
\item \textbf{Curve above then below:} Left skewness
\item \textbf{Points at ends deviate:} Extreme values
\end{itemize}

\vspace{0.3cm}
\textbf{R code:}
\begin{verbatim}
qqnorm(residuals(model))
qqline(residuals(model))
\end{verbatim}

\vspace{0.2cm}
\textbf{Note:} Minor departures are okay, especially with small samples
\end{frame}

\begin{frame}[fragile]
\frametitle{Scale-Location Plot}
\textbf{Purpose:} Check homoscedasticity (constant variance)

\vspace{0.3cm}
\textbf{What it shows:} Square root of standardized residuals vs fitted values

\vspace{0.3cm}
\textbf{What to look for:}
\begin{itemize}
\item \textbf{Good:} Random scatter with horizontal trend line
\item \textbf{Bad:} Upward or downward trend
\end{itemize}

\vspace{0.5cm}
\textbf{Interpretation:}
\begin{itemize}
\item Horizontal line $\Rightarrow$ constant variance
\item Upward slope $\Rightarrow$ variance increases with fitted values
\item Downward slope $\Rightarrow$ variance decreases with fitted values
\end{itemize}

\vspace{0.3cm}
\textbf{R code:}
\begin{verbatim}
plot(model, which=3)
\end{verbatim}
\end{frame}

\begin{frame}[fragile]
\frametitle{Residuals vs Leverage Plot}
\textbf{Purpose:} Identify influential observations

\vspace{0.3cm}
\textbf{Key Concepts:}
\begin{itemize}
\item \textbf{Leverage:} How far $x_i$ is from $\bar{x}$ (potential influence)
\item \textbf{Cook's Distance:} Actual influence on fitted values
\item \textbf{Influential point:} High leverage AND large residual
\end{itemize}

\vspace{0.5cm}
\textbf{What to look for:}
\begin{itemize}
\item Points outside Cook's distance contours (dashed lines)
\item Cook's distance $> 0.5$ or $> 1$ suggests influential observation
\end{itemize}

\vspace{0.3cm}
\textbf{R code:}
\begin{verbatim}
plot(model, which=5)
\end{verbatim}

\vspace{0.2cm}
\textbf{Action:} Investigate influential points, but don't automatically remove them
\end{frame}

\begin{frame}[fragile]
\frametitle{Creating All Diagnostic Plots in R}
\textbf{Method 1: All four plots at once}
\small
\begin{verbatim}
par(mfrow=c(2,2))
plot(model)
par(mfrow=c(1,1))
\end{verbatim}

\vspace{0.3cm}
\normalsize
\textbf{Method 2: Individual plots}
\small
\begin{verbatim}
# 1. Residuals vs Fitted
plot(model, which=1)

# 2. Q-Q plot
plot(model, which=2)

# 3. Scale-Location
plot(model, which=3)

# 4. Residuals vs Leverage
plot(model, which=5)
\end{verbatim}
\end{frame}

\section{Outliers and Influential Points}

\begin{frame}
\frametitle{Types of Unusual Observations}
\textbf{Three types to distinguish:}

\vspace{0.3cm}
\textbf{1. Outlier:}
\begin{itemize}
\item Large residual (far from regression line)
\item May or may not be influential
\item Check for data entry errors
\end{itemize}

\vspace{0.3cm}
\textbf{2. High leverage point:}
\begin{itemize}
\item Extreme $x$ value (far from $\bar{x}$)
\item Potential to be influential
\item May not have large residual
\end{itemize}

\vspace{0.3cm}
\textbf{3. Influential point:}
\begin{itemize}
\item High leverage AND large residual
\item Removing it substantially changes regression line
\item Check Cook's distance
\end{itemize}
\end{frame}

\begin{frame}
\frametitle{Dealing with Influential Points}
\textbf{DO NOT automatically remove outliers!}

\vspace{0.3cm}
\textbf{Steps to take:}
\begin{enumerate}
\item \textbf{Check for errors:}
  \begin{itemize}
  \item Data entry mistakes
  \item Measurement errors
  \item Wrong units
  \end{itemize}

\item \textbf{Investigate context:}
  \begin{itemize}
  \item Is it a legitimate observation?
  \item Does it represent a special subgroup?
  \item Was there a protocol violation?
  \end{itemize}

\item \textbf{Perform sensitivity analysis:}
  \begin{itemize}
  \item Fit model with and without the point
  \item Report both analyses
  \item Discuss impact on conclusions
  \end{itemize}

\item \textbf{Consider robust methods:}
  \begin{itemize}
  \item Robust regression (less sensitive to outliers)
  \item Transformations
  \end{itemize}
\end{enumerate}
\end{frame}

\begin{frame}[fragile]
\frametitle{Example: Identifying Influential Points in R}
\small
\begin{verbatim}
# Fit model
model <- lm(y ~ x, data=mydata)

# Calculate Cook's distance
cooks <- cooks.distance(model)

# Identify influential points (Cook's D > 0.5)
influential <- which(cooks > 0.5)
print(influential)

# Show influential observations
mydata[influential, ]

# Compare models with/without influential points
model_full <- lm(y ~ x, data=mydata)
model_reduced <- lm(y ~ x, data=mydata[-influential,])

summary(model_full)
summary(model_reduced)
\end{verbatim}
\end{frame}

\section{Transformations}

\begin{frame}
\frametitle{When to Transform Data}
\textbf{Transformations can help when:}
\begin{itemize}
\item Relationship is non-linear
\item Variance is non-constant
\item Residuals are non-normal
\item Interpretation requires it (e.g., multiplicative effects)
\end{itemize}

\vspace{0.5cm}
\textbf{Common transformations:}
\begin{itemize}
\item \textbf{Log:} $\log(Y)$ or $\log(X)$ -- for right skewness, exponential growth
\item \textbf{Square root:} $\sqrt{Y}$ -- for count data, mild right skewness
\item \textbf{Reciprocal:} $1/Y$ or $1/X$ -- for rates, strong right skewness
\item \textbf{Square:} $Y^2$ or $X^2$ -- for quadratic relationships
\end{itemize}

\vspace{0.3cm}
\textbf{Box-Cox transformation:} Data-driven choice of power transformation
\end{frame}

\begin{frame}
\frametitle{Log Transformation}
\textbf{Most common transformation in biostatistics}

\vspace{0.3cm}
\textbf{Model forms:}
\begin{enumerate}
\item \textbf{Log-linear:} $\log(Y) = \beta_0 + \beta_1 X$
  \begin{itemize}
  \item Interpretation: 1-unit increase in $X$ $\Rightarrow$ $(e^{\beta_1} - 1) \times 100\%$ change in $Y$
  \end{itemize}

\item \textbf{Linear-log:} $Y = \beta_0 + \beta_1 \log(X)$
  \begin{itemize}
  \item Interpretation: 1\% increase in $X$ $\Rightarrow$ $\beta_1/100$ unit change in $Y$
  \end{itemize}

\item \textbf{Log-log:} $\log(Y) = \beta_0 + \beta_1 \log(X)$
  \begin{itemize}
  \item Interpretation: $\beta_1$ is the elasticity (1\% change in $X$ $\Rightarrow$ $\beta_1\%$ change in $Y$)
  \end{itemize}
\end{enumerate}

\vspace{0.3cm}
\textbf{When to use:}
\begin{itemize}
\item Right-skewed data
\item Multiplicative relationships
\item Relative changes more meaningful than absolute
\end{itemize}
\end{frame}

\begin{frame}[fragile]
\frametitle{Example: Log Transformation in R}
\textbf{Scenario:} BMI vs medical cost (right-skewed)

\small
\begin{verbatim}
# Original model (may violate assumptions)
model_original <- lm(cost ~ bmi, data=health)

# Log-transformed response
model_log <- lm(log(cost) ~ bmi, data=health)

# Compare diagnostics
par(mfrow=c(2,2))
plot(model_original)
plot(model_log)
par(mfrow=c(1,1))

# Interpretation
summary(model_log)
# If beta1 = 0.05, then 1-unit increase in BMI
# increases cost by (exp(0.05)-1)*100 = 5.1%
\end{verbatim}
\end{frame}

\begin{frame}
\frametitle{Polynomial Regression}
\textbf{For non-linear relationships with turning points}

\vspace{0.3cm}
\textbf{Quadratic model:}
$$Y = \beta_0 + \beta_1 X + \beta_2 X^2 + \epsilon$$

\vspace{0.3cm}
\textbf{When to use:}
\begin{itemize}
\item U-shaped or inverted U-shaped relationship
\item Single peak or valley
\item Example: age vs bone density (increases then decreases)
\end{itemize}

\vspace{0.5cm}
\textbf{Interpretation:}
\begin{itemize}
\item $\beta_1$: slope at $X = 0$
\item $\beta_2$: rate of change in slope
\item If $\beta_2 > 0$: U-shaped (convex)
\item If $\beta_2 < 0$: inverted U-shaped (concave)
\end{itemize}

\vspace{0.3cm}
\textbf{Warning:} Don't extrapolate beyond data range!
\end{frame}

\begin{frame}[fragile]
\frametitle{Polynomial Regression in R}
\small
\begin{verbatim}
# Create quadratic term
mydata$age_squared <- mydata$age^2

# Fit quadratic model
model_quad <- lm(bone_density ~ age + age_squared,
                 data=mydata)

# Alternative: use poly() function
model_quad2 <- lm(bone_density ~ poly(age, 2, raw=TRUE),
                  data=mydata)

summary(model_quad)

# Find peak (maximum bone density)
# Peak at: -beta1 / (2 * beta2)
beta1 <- coef(model_quad)[2]
beta2 <- coef(model_quad)[3]
peak_age <- -beta1 / (2 * beta2)
print(paste("Peak at age:", round(peak_age, 1)))
\end{verbatim}
\end{frame}

\section{Hypothesis Testing and Confidence Intervals}

\begin{frame}
\frametitle{Testing the Slope}
\textbf{Null hypothesis:} $H_0: \beta_1 = 0$ (no linear relationship)

\textbf{Alternative:} $H_A: \beta_1 \neq 0$ (linear relationship exists)

\vspace{0.5cm}
\textbf{Test statistic:}
$$t = \frac{b_1 - 0}{SE(b_1)} \sim t_{n-2}$$

where $SE(b_1) = \frac{s_e}{\sqrt{\sum(x_i - \bar{x})^2}}$ and $s_e = \sqrt{\frac{SSE}{n-2}}$

\vspace{0.5cm}
\textbf{Decision:}
\begin{itemize}
\item Reject $H_0$ if $|t| > t_{\alpha/2, n-2}$ or $p < \alpha$
\item Fail to reject if $|t| \leq t_{\alpha/2, n-2}$ or $p \geq \alpha$
\end{itemize}

\vspace{0.3cm}
\textbf{Note:} This tests whether $X$ is useful for predicting $Y$
\end{frame}

\begin{frame}
\frametitle{Confidence Interval for Slope}
\textbf{95\% CI for $\beta_1$:}
$$b_1 \pm t_{0.975, n-2} \times SE(b_1)$$

\vspace{0.5cm}
\textbf{Interpretation:}
\begin{itemize}
\item We are 95\% confident the true slope is in this interval
\item If interval excludes 0, reject $H_0: \beta_1 = 0$
\item Width reflects precision of estimate
\end{itemize}

\vspace{0.5cm}
\textbf{Example:}
\begin{itemize}
\item $b_1 = 2.5$ (estimated slope)
\item $SE(b_1) = 0.8$
\item $n = 30$, so $t_{0.975, 28} = 2.048$
\item 95\% CI: $2.5 \pm 2.048(0.8) = (0.86, 4.14)$
\item Interpretation: For each 1-unit increase in $X$, $Y$ increases between 0.86 and 4.14 units
\end{itemize}
\end{frame}

\begin{frame}
\frametitle{Confidence vs Prediction Intervals}
\textbf{Two types of intervals for predictions:}

\vspace{0.3cm}
\textbf{1. Confidence Interval for Mean Response:}
\begin{itemize}
\item Estimates the \textbf{average} $Y$ for a given $X = x_0$
\item Answers: What is the mean outcome for this X?
\item Narrower interval
\item Example: Mean blood pressure for all 50-year-olds
\end{itemize}

\vspace{0.3cm}
\textbf{2. Prediction Interval for Individual Response:}
\begin{itemize}
\item Predicts a \textbf{single} $Y$ for a given $X = x_0$
\item Answers: What will this individual's outcome be?
\item Wider interval (includes individual variability)
\item Example: Blood pressure for one specific 50-year-old
\end{itemize}

\vspace{0.3cm}
\textbf{Key difference:} Prediction interval accounts for both:
\begin{enumerate}
\item Uncertainty in estimating the mean (same as CI)
\item Variability of individuals around the mean (extra)
\end{enumerate}
\end{frame}

\begin{frame}[fragile]
\frametitle{Computing Intervals in R}
\small
\begin{verbatim}
# Fit model
model <- lm(sbp ~ age, data=health)

# New data point
new_data <- data.frame(age = 50)

# Confidence interval for mean SBP at age 50
predict(model, newdata=new_data,
        interval="confidence", level=0.95)
# Example output: fit=130, lwr=128, upr=132

# Prediction interval for individual at age 50
predict(model, newdata=new_data,
        interval="prediction", level=0.95)
# Example output: fit=130, lwr=110, upr=150

# Note: Prediction interval is wider!
\end{verbatim}
\end{frame}

\section{Presenting Statistical Results}

\begin{frame}
\frametitle{Elements of a Good Statistical Presentation}
\textbf{For your final projects, include:}

\vspace{0.3cm}
\begin{enumerate}
\item \textbf{Clear research question}
  \begin{itemize}
  \item What are you trying to learn?
  \end{itemize}

\item \textbf{Descriptive statistics}
  \begin{itemize}
  \item Sample size, means, SDs
  \item Summary tables
  \end{itemize}

\item \textbf{Visualizations}
  \begin{itemize}
  \item Scatterplots, histograms, boxplots
  \item Clear labels and legends
  \end{itemize}

\item \textbf{Statistical methods}
  \begin{itemize}
  \item Which tests/models did you use?
  \item Why did you choose them?
  \end{itemize}

\item \textbf{Results}
  \begin{itemize}
  \item Effect sizes, confidence intervals, p-values
  \item Interpretation in context
  \end{itemize}

\item \textbf{Limitations and assumptions}
  \begin{itemize}
  \item What could go wrong?
  \end{itemize}
\end{enumerate}
\end{frame}

\begin{frame}
\frametitle{Reporting Regression Results}
\textbf{Essential components:}

\vspace{0.3cm}
\begin{itemize}
\item \textbf{Sample size:} $n = 250$ observations
\item \textbf{Model equation:} $\hat{Y} = b_0 + b_1 X$
\item \textbf{Coefficient estimates:}
  \begin{itemize}
  \item Point estimates
  \item Standard errors
  \item 95\% confidence intervals
  \end{itemize}
\item \textbf{Statistical significance:} t-statistics, p-values
\item \textbf{Model fit:} $R^2$, residual standard error
\item \textbf{Interpretation in context:}
  \begin{itemize}
  \item What does the slope mean practically?
  \item Clinical/scientific significance vs statistical
  \end{itemize}
\item \textbf{Diagnostics:} Did assumptions hold?
\end{itemize}

\vspace{0.3cm}
\textbf{Example:} ``For each 1-unit increase in BMI, systolic BP increased by 1.2 mmHg (95\% CI: 0.8-1.6, p<0.001). The model explained 23\% of the variance in SBP ($R^2=0.23$).''
\end{frame}

\begin{frame}
\frametitle{Common Mistakes to Avoid}
\textbf{Don't:}
\begin{itemize}
\item Confuse correlation with causation
  \begin{itemize}
  \item ``X \textbf{is associated with} Y'' not ``X \textbf{causes} Y''
  \end{itemize}
\item Over-interpret R-squared
  \begin{itemize}
  \item Low $R^2$ doesn't mean the relationship isn't important
  \end{itemize}
\item Extrapolate beyond the data range
  \begin{itemize}
  \item Don't predict for X values outside your sample
  \end{itemize}
\item Ignore influential points without investigation
\item Report only p-values without effect sizes
  \begin{itemize}
  \item \textbf{Always} report confidence intervals!
  \end{itemize}
\item Forget to check assumptions
  \begin{itemize}
  \item Always create diagnostic plots
  \end{itemize}
\item Use jargon without explanation
  \begin{itemize}
  \item Explain terms for non-statistical audience
  \end{itemize}
\end{itemize}
\end{frame}

\begin{frame}
\frametitle{Creating Effective Visualizations}
\textbf{Best practices:}

\vspace{0.3cm}
\begin{itemize}
\item \textbf{Always label axes} with units
\item \textbf{Include informative titles}
\item \textbf{Use appropriate plot types}
  \begin{itemize}
  \item Scatterplot for continuous X and Y
  \item Boxplot for categorical X, continuous Y
  \end{itemize}
\item \textbf{Show the data and the model}
  \begin{itemize}
  \item Points + regression line
  \item Confidence/prediction bands
  \end{itemize}
\item \textbf{Use color wisely}
  \begin{itemize}
  \item Colorblind-friendly palettes
  \item Sufficient contrast
  \end{itemize}
\item \textbf{Keep it simple}
  \begin{itemize}
  \item One main message per plot
  \item Remove unnecessary elements
  \end{itemize}
\end{itemize}
\end{frame}

\begin{frame}[fragile]
\frametitle{Example: Publication-Quality Plot}
\small
\begin{verbatim}
# Create professional scatterplot with regression
model <- lm(sbp ~ age, data=health)

# Set up plot
plot(health$age, health$sbp,
     xlab="Age (years)",
     ylab="Systolic Blood Pressure (mmHg)",
     main="Relationship between Age and Blood Pressure",
     pch=16, col=rgb(0,0,0,0.3), cex=1.2)

# Add regression line
abline(model, col="red", lwd=2)

# Add confidence band
new_age <- seq(min(health$age), max(health$age), 0.1)
pred_ci <- predict(model,
                   newdata=data.frame(age=new_age),
                   interval="confidence")
lines(new_age, pred_ci[,"lwr"], col="blue", lty=2)
lines(new_age, pred_ci[,"upr"], col="blue", lty=2)

# Add legend
legend("topleft",
       c("Observed", "Fitted line", "95% CI"),
       col=c("black", "red", "blue"),
       lty=c(NA, 1, 2), pch=c(16, NA, NA))
\end{verbatim}
\end{frame}

\section{Case Study: Complete Analysis}

\begin{frame}
\frametitle{Case Study: FEV and Age}
\textbf{Research Question:} How does lung function (FEV) change with age in children?

\vspace{0.3cm}
\textbf{Data:} 654 children aged 3-19 years
\begin{itemize}
\item \textbf{Outcome:} FEV (forced expiratory volume, liters)
\item \textbf{Predictor:} Age (years)
\end{itemize}

\vspace{0.3cm}
\textbf{Hypothesis:} FEV increases with age as lungs develop

\vspace{0.3cm}
\textbf{Analysis Plan:}
\begin{enumerate}
\item Descriptive statistics and visualization
\item Fit simple linear regression
\item Check assumptions
\item Interpret results
\item Consider alternatives if needed
\end{enumerate}
\end{frame}

\begin{frame}[fragile]
\frametitle{Step 1: Descriptive Statistics}
\small
\begin{verbatim}
# Summary statistics
summary(fev_data$age)
#   Min. 1st Qu.  Median    Mean 3rd Qu.    Max.
#   3.00    8.00   10.00    9.93   12.00   19.00

summary(fev_data$fev)
#   Min. 1st Qu.  Median    Mean 3rd Qu.    Max.
#  0.791   1.981   2.547   2.637   3.119   5.793

sd(fev_data$age)      # 2.95 years
sd(fev_data$fev)      # 0.87 liters

# Correlation
cor(fev_data$age, fev_data$fev)
# r = 0.877 (strong positive correlation)
\end{verbatim}
\end{frame}

\begin{frame}[fragile]
\frametitle{Step 2: Initial Visualization}
\small
\begin{verbatim}
# Scatterplot
plot(fev_data$age, fev_data$fev,
     xlab="Age (years)",
     ylab="FEV (liters)",
     main="FEV vs Age in Children",
     pch=16, col=rgb(0,0,1,0.3))

# Observations:
# - Strong positive relationship
# - Possible non-linearity (steeper for younger)
# - Greater variability for older children
\end{verbatim}

\vspace{0.2cm}
\normalsize
\textbf{Initial impression:} Linear model may work, but check diagnostics
\end{frame}

\begin{frame}[fragile]
\frametitle{Step 3: Fit Linear Model}
\small
\begin{verbatim}
# Fit simple linear regression
model <- lm(fev ~ age, data=fev_data)
summary(model)

# Coefficients:
#             Estimate Std. Error t value Pr(>|t|)
# (Intercept)  0.43161    0.06328   6.821  1.8e-11 ***
# age          0.22204    0.00608  36.515  < 2e-16 ***
# ---
# Residual standard error: 0.358 on 652 degrees of freedom
# Multiple R-squared:  0.6714,    Adjusted R-squared:  0.6709
# F-statistic: 1333 on 1 and 652 DF,  p-value: < 2.2e-16

# Interpretation:
# FEV increases by 0.222 liters per year of age
# Model explains 67.1% of variance
\end{verbatim}
\end{frame}

\begin{frame}[fragile]
\frametitle{Step 4: Check Assumptions}
\small
\begin{verbatim}
# Diagnostic plots
par(mfrow=c(2,2))
plot(model)
par(mfrow=c(1,1))

# Observations:
# 1. Residuals vs Fitted: slight curve (non-linearity?)
# 2. Q-Q plot: reasonably linear (normality OK)
# 3. Scale-Location: some funneling (heteroscedasticity?)
# 4. Residuals vs Leverage: a few high leverage points

# Potential issues:
# - Non-linearity (could try quadratic or log)
# - Non-constant variance (more spread for older ages)
\end{verbatim}
\end{frame}

\begin{frame}[fragile]
\frametitle{Step 5: Consider Alternatives}
\textbf{Given diagnostic concerns, try quadratic model:}

\small
\begin{verbatim}
# Quadratic model
model_quad <- lm(fev ~ age + I(age^2), data=fev_data)
summary(model_quad)

# Coefficients:
#             Estimate Std. Error t value Pr(>|t|)
# (Intercept)  1.04423    0.10055   10.39  < 2e-16 ***
# age         -0.01626    0.02273   -0.72    0.475
# I(age^2)     0.01201    0.00109   11.00  < 2e-16 ***
# ---
# Multiple R-squared:  0.7198,    Adjusted R-squared:  0.7189

# Better R-squared (72% vs 67%)
# Check diagnostics again...
\end{verbatim}
\end{frame}

\begin{frame}
\frametitle{Step 6: Interpret and Report}
\textbf{Final Model:} Quadratic relationship

$$\text{FEV} = 1.04 - 0.016 \times \text{Age} + 0.012 \times \text{Age}^2$$

\vspace{0.3cm}
\textbf{Key Findings:}
\begin{itemize}
\item FEV increases with age in children, but non-linearly
\item Rate of increase accelerates with age
\item Model explains 72\% of variance
\item All coefficients highly significant ($p < 0.001$)
\end{itemize}

\vspace{0.3cm}
\textbf{Conclusion:}
\begin{itemize}
\item Strong positive relationship between age and lung function
\item Quadratic model fits better than linear
\item Results consistent with biological lung development
\item Caution: Don't extrapolate to adults (age > 19)
\end{itemize}
\end{frame}

\section{Final Project Guidance}

\begin{frame}
\frametitle{Final Project Presentations}
\textbf{Format:}
\begin{itemize}
\item 10-12 minute presentation
\item 3-5 minutes for questions
\item PowerPoint, Beamer, or Google Slides
\end{itemize}

\vspace{0.3cm}
\textbf{Required content:}
\begin{enumerate}
\item Introduction and research question
\item Data description and summary statistics
\item Statistical methods used
\item Results (tables, figures)
\item Interpretation and conclusions
\item Limitations
\end{enumerate}

\vspace{0.3cm}
\textbf{Grading criteria:}
\begin{itemize}
\item Clarity and organization (20\%)
\item Appropriate statistical methods (30\%)
\item Correct interpretation (25\%)
\item Quality of visualizations (15\%)
\item Presentation delivery (10\%)
\end{itemize}
\end{frame}

\begin{frame}
\frametitle{Tips for Successful Presentations}
\textbf{Do:}
\begin{itemize}
\item Practice your presentation multiple times
\item Speak clearly and at moderate pace
\item Make eye contact with audience
\item Explain statistical terms
\item Focus on interpretation, not just methods
\item Use clear, uncluttered slides
\item Tell a story with your data
\end{itemize}

\vspace{0.3cm}
\textbf{Don't:}
\begin{itemize}
\item Read directly from slides
\item Use too much text on slides
\item Skip over assumptions or limitations
\item Ignore questions you don't know -- say ``good question, I'll look into that''
\item Go over time limit
\end{itemize}
\end{frame}

\begin{frame}
\frametitle{Common Presentation Pitfalls}
\textbf{Avoid these mistakes:}

\vspace{0.3cm}
\begin{enumerate}
\item \textbf{Too much technical detail}
  \begin{itemize}
  \item Focus on interpretation, not formulas
  \end{itemize}

\item \textbf{No context for numbers}
  \begin{itemize}
  \item Always interpret in practical terms
  \end{itemize}

\item \textbf{Cluttered visualizations}
  \begin{itemize}
  \item One clear message per plot
  \end{itemize}

\item \textbf{Ignoring limitations}
  \begin{itemize}
  \item Acknowledge what your analysis can't tell you
  \end{itemize}

\item \textbf{Poor time management}
  \begin{itemize}
  \item Practice with a timer
  \end{itemize}

\item \textbf{Not checking assumptions}
  \begin{itemize}
  \item Always show diagnostic plots
  \end{itemize}
\end{enumerate}
\end{frame}

\begin{frame}
\frametitle{Summary}
\textbf{This week we covered:}
\begin{itemize}
\item \textbf{Regression diagnostics:}
  \begin{itemize}
  \item Four key assumptions (LINE)
  \item Diagnostic plots
  \item Identifying influential points
  \end{itemize}

\item \textbf{Transformations:}
  \begin{itemize}
  \item Log, square root, polynomial
  \item When and why to transform
  \end{itemize}

\item \textbf{Inference:}
  \begin{itemize}
  \item Hypothesis tests for slope
  \item Confidence and prediction intervals
  \end{itemize}

\item \textbf{Communication:}
  \begin{itemize}
  \item Reporting results clearly
  \item Creating effective visualizations
  \item Using AI to enhance presentations
  \end{itemize}
\end{itemize}

\vspace{0.5cm}
\textbf{Final Project Presentations:} Next week!

\vspace{0.3cm}
\textbf{Prepare:} Practice your presentation, check all analyses, create clear visuals
\end{frame}

\begin{frame}
\frametitle{Questions?}
\begin{center}
\Large{Thank you!}

\vspace{0.5cm}

\textbf{Good luck with your final projects!}

\vspace{1cm}

\textbf{Email:} john.molitor@oregonstate.edu\\
\textbf{Office:} 153 Milam
\end{center}
\end{frame}

\end{document}
